\section{Introduction: Relational Complexity and Conserved Information}\label{sec:introduction}
\begin{principle}[title=Quantum Relativity]
Quantum Relativity is a theory of relational complexity governed by conservation of information. In compact form, QR studies relational complexity under information-preserving transduction.
\end{principle}
The subject is the organization of registered relations and the transformations that preserve a complete account of them. A registered object $R$ is the thing whose relational structure is under consideration. A complexity functional $\mathcal C(R)$ asks how much such structure is present; phase, Mode, and other readouts describe aspects of its organization. The transducer is the mechanism of change, while complete information preservation is the admissibility constraint. This thesis defines the program. It does not by itself select a unique complexity measure or establish a physical conservation theorem.

The primitive question is operational: what must occur for one well-defined relation to become another while retaining the information required to account for the transition? A local transducer receives an admissible input, produces a propagated component and a complementary retained output, and completes a handoff to an adjacent address:
% Original equation number(s) 1; expression unchanged.
\setcounter{equation}{0}
\begin{equation}
\begin{aligned}
\text{relation}&\longrightarrow\text{admissible transduction}\\
&\longrightarrow\text{propagated component + retained receipt}\\
&\longrightarrow\text{adjacent handoff}\longrightarrow\text{history}.
\end{aligned}
\end{equation}

The joint output is the object of information conservation. A completed handoff is the primitive propagation event; composition supplies its historical placement. The scalar encoding $S(R)=G^{\mathcal C(R)}$ is a readout of the relational object, not a second substance moving through this chain. Its relay requires an admitted update and enough local context to determine the new value.

Quantum theory organizes preparations, observables, and outcome probabilities; general relativity organizes matter and causal relations through a dynamical metric \citep{watrous2018,carroll1997}. QR investigates whether compatible descriptions of both can be obtained from the same information-preserving relational process. The Thalean carrier $G_{60}\cong\mathrm{AT4val}[60,6]$ provides a finite setting in which quotient distinctions, receipts, phase actions, and histories can be stated exactly. The architecture draws on reversible computation, relational quantum theory, quantum cellular automata, and causal-order programs \citep{bennett1973,rovelli1996,arrighi2019,sorkin2003}. Its task is to derive compatible operations and observables from declared structures, rather than attach independent quantum and geometric models to a common graph.

\subsection*{Reference without privilege}
The motivating formulation is methodological: ``GR says no frame is privileged. QR says: agreed. Select one. Then we can start.'' This is not a quotation from a relativistic theorem or a claim that the finite carrier derives relativistic frame transformations. It distinguishes a reference selected for an execution from a frame declared physically preferred by the laws.

A reader reconstructing a calculation and an apparatus retaining an experimental record both work through specified interfaces. The account accumulates at that registration, with its retained context. Selection makes local coordinates definite; incidence supplies the place at which one account can be compared with another. Orientation reconciliation then requires a compatible comparison law, not an absolute sign assigned to the entire universe.

Sections~\ref{sec:executingreference} and \ref{sec:incidencereconciliation} formalize this distinction. The new criterion states exactly when two lifted sign actions over a common incidence domain admit a covariant comparison. It does not assume that the required native history--analyzer domain has already been constructed. A successful comparison may preserve simultaneous reversal of both references; it need not select a privileged hand, center, or physical observer.

\subsection*{Naming and intellectual background}\label{sec:naming}
\textbf{Thalean} refers to Thales of Miletus. Aristotle reports water as Thales' originating principle, or \emph{arch\={e}}; ``everything is water'' is used here as a traditional shorthand, not a surviving verbatim quotation by Thales \citep[Book~I.3, 983b6--27]{aristotleMetaphysics}. The name evokes the turn associated with early Milesian inquiry toward explanation through nature's own principles rather than external divine intervention \citep[Sections~1--2]{curdPresocratics}. The intended \emph{abandonment of caprice} is methodological: seek a lawful account instead of an arbitrary exception. It is not a claim that Thales rejected every conception of divinity, nor does QR adopt water as a literal physical substrate. In QR the corresponding search concerns relational structure and the information-preserving transformations beneath differing appearances.

\textbf{Thalion} has a separate, serendipitous literary resonance. Tolkien's epithet \emph{H\'{u}rin Thalion} is rendered ``H\'{u}rin the Steadfast''; the associated sense also includes strength \citep{fisherThalion}. Here the coincidence is a mnemonic, not mathematical provenance. A Thalion is intended to be \emph{informationally steadfast}: its relational organization may change, but admissible transformation does not abandon the complete account required for reconstruction.

\subsection*{Route through the paper}
Part~I supplies the carrier, cover family and relational kernel. Parts~II--IV introduce registered objects, selected executing reference, complete transduction and conservation. Part~V develops scalar relay, phase, Mode and incidence-based orientation reconciliation. Part~VI treats history, literal phase transport and Thalean time. Parts~VII--VIII retain causal support and measure, then distinguish signed analyzer presentations from native instrument actions. Part~IX states the quantum and geometric reductions and their common-domain requirements. Part~X contains the conservative-generator, gyroscopic and Maxwell comparison models. Part~XI gives the open correspondence program, followed by tests, guardrails and synthesis.

Inherited section and equation identifiers remain fixed for audit continuity. Appendix~H is the proof-status ledger; Appendices~M and N retain the scalar-relay and conservative-calculation contracts. Appendix~O records the new source integrity, incidence and signed-group replays, and the qualification of the stronger Audit018K history-state interpretation. The aim is one continuous argument, with definitions, established finite implications and missing physical maps kept distinct.
